\section{True projective residue valuations from original relevant graded seminorms}

\begin{theorem}[A relevant graded seminorm defines a true native Proj valued point]
\label{thm:fc-graded-native-valued-point}
\leanok
\lean{EnvelopeProperty.Berkovich.RelevantGradedSpectrum.localValuation, EnvelopeProperty.Berkovich.RelevantGradedSpectrum.localValuation_supp, EnvelopeProperty.Berkovich.RelevantGradedSpectrum.homogeneousLocalValuation, EnvelopeProperty.Berkovich.RelevantGradedSpectrum.homogeneousLocalValuation_supp, EnvelopeProperty.Berkovich.RelevantGradedSpectrum.nativeStalkValuation, EnvelopeProperty.Berkovich.RelevantGradedSpectrum.nativeStalkValuation_supp, EnvelopeProperty.Berkovich.RelevantGradedSpectrum.nativePoint, EnvelopeProperty.Berkovich.RelevantGradedSpectrum.nativePoint_point, EnvelopeProperty.Berkovich.RelevantGradedSpectrum.nativePoint_residue}
\uses{lem:fc-native-model-total-grading}
For an actual graded commutative ring, every original relevant graded
real seminorm defines a canonical native valued point of its actual
Proj. Its scheme support is the SAME original homogeneous prime kernel.
Its projective residue-field valuation retains the original valuation
of every actual projective stalk residue.
\end{theorem}
\begin{proof}
Extend the original valuation to the ordinary local ring at its kernel.
The valuation of a fraction vanishes exactly when its numerator is in
that original kernel, so the support is precisely the actual local
maximal ideal. Restrict this valuation to the TRUE homogeneous
degree-zero local ring. The true homogeneous localization unit criterion
shows that its support is again exactly the actual local maximal ideal.
Transport along the actual native Proj stalk isomorphism and descend
through its true residue quotient. This gives the stated valued point
and retains the whole original homogeneous support and every genuine
stalk residue. The generally larger affine-cone residue field is never
identified with the projective residue field.
\end{proof}

\begin{theorem}[Actual native homogeneous fraction sections keep their original values]
\label{thm:fc-graded-native-valued-fractions}
\leanok
\lean{EnvelopeProperty.Berkovich.RelevantGradedSpectrum.homogeneousLocalValuation_mapId_chart, EnvelopeProperty.Berkovich.RelevantGradedSpectrum.nativePoint_eval_awayToSection, EnvelopeProperty.Berkovich.RelevantGradedSpectrum.nativePoint_eval_mk}
\uses{thm:fc-graded-native-valued-point}
Every true native homogeneous chart section is evaluated by the original
seminorm extended to that SAME principal localization. In particular,
an original numerator $p$ divided by an original homogeneous denominator
$f^d$ has exactly the original value $|p|/|f|^d$ whenever $|f|\ne0$.
\end{theorem}
\begin{proof}
Every true homogeneous chart element is an actual same-degree fraction.
The native fraction section's germ is its actual homogeneous-local-ring
image followed by the inverse of the true projective stalk isomorphism.
The residue construction retains that stalk value. The actual ordinary
localized valuation evaluates that SAME fraction as the original
numerator value times the inverse denominator value; multiplicativity
gives the displayed power formula. This applies to all actual original
fractions, including zero numerators.
\end{proof}
